% 멀티카메라 캘리브레이션 설명서 1~11장 발표자료
%
% 컴파일: XeLaTeX 또는 LuaLaTeX 로 컴파일해야 한글이 나온다.
%   Overleaf: 왼쪽 위 Menu > Compiler > XeLaTeX 선택
%   로컬:     xelatex calibration_1_11.tex   (두 번 실행)
%
% 원문: CALIBRATION_EXPLANATION_LATEX.md 의 1장부터 11장까지

\documentclass[aspectratio=169,11pt]{beamer}

\usepackage{kotex}
\usepackage{amsmath}
\usepackage{amssymb}
\usepackage{booktabs}
\usepackage{array}
\usepackage{tabularx}

\usetheme{Madrid}
\usecolortheme{whale}
\setbeamertemplate{navigation symbols}{}
\setbeamertemplate{caption}[numbered]
\setbeamerfont{frametitle}{size=\large}
\setbeamerfont{block title}{size=\normalsize}

% 표 안에서 줄 간격을 조금 넓힌다
\renewcommand{\arraystretch}{1.3}

\title[멀티카메라 캘리브레이션]{멀티카메라 캘리브레이션 완전 설명서}
\subtitle{1장 -- 11장: 좌표계부터 세 가지 FK 방식까지}
\author{}
\institute{}
\date{}

\begin{document}

\begin{frame}
  \titlepage
\end{frame}

\begin{frame}{다루는 내용}
  \tableofcontents
\end{frame}

%==================================================================
\section{1. 왜 카메라마다 큐브의 좌표가 다를까?}
%==================================================================

\begin{frame}{1. 왜 카메라마다 큐브의 좌표가 다를까?}
  카메라마다 자기 자신을 원점으로 사용하기 때문이다.
  \begin{itemize}
    \item 카메라 $0$은 큐브가 자기 앞에 있다고 말한다.
    \item 카메라 $1$은 같은 큐브가 자기 오른쪽에 있다고 말한다.
    \item 그리퍼 카메라는 같은 큐브가 자기 아래에 있다고 말할 수 있다.
  \end{itemize}
  \vfill
  세 답은 서로 모순이 아니다. 기준 좌표계가 다를 뿐이다.
  \vfill
  \begin{block}{그래서 해야 할 일}
    모든 관측을 공통 기준인 로봇 베이스 좌표계 $B$로 옮긴다.
  \end{block}
\end{frame}

%==================================================================
\section{2. 좌표계 기호}
%==================================================================

\begin{frame}{2. 먼저 좌표계 기호를 정리한다}
  \begin{center}
    \begin{tabular}{cl}
      \toprule
      기호 & 의미 \\
      \midrule
      $B$      & 로봇 베이스 좌표계 \\
      $C_i$    & $i$번째 고정 카메라 좌표계 \\
      $G$      & 로봇 그리퍼 좌표계 \\
      $C_g$    & 그리퍼에 부착된 카메라 좌표계 \\
      $O$      & 큐브 또는 캘리브레이션 물체 좌표계 \\
      $s$      & 큐브를 놓은 세트 번호 \\
      $e$      & 그리퍼 카메라 촬영 이벤트 번호 \\
      $s(e)$   & 이벤트 $e$가 속한 세트 번호 \\
      \bottomrule
    \end{tabular}
  \end{center}
\end{frame}

\begin{frame}{2. 변환 표기를 읽는 방법}
  이 문서에서는 다음과 같은 변환 표기를 사용한다.
  \[
    {}^{A}\mathbf{T}_{B}
  \]
  \begin{block}{읽는 방법}
    좌표계 $B$에서 표현된 값을 좌표계 $A$에서 표현하도록 바꾸는 변환이다.
  \end{block}
  \vfill
  위 첨자 $A$는 \textbf{도착} 좌표계이고, 아래 첨자 $B$는 \textbf{출발} 좌표계다.
\end{frame}

%==================================================================
\section{3. 변환행렬}
%==================================================================

\begin{frame}{3. 변환행렬은 무엇인가?}
  하나의 강체 자세는 회전과 이동으로 이루어진 $4 \times 4$ 동차변환행렬로 표현한다.
  \[
    \mathbf{T}
    =
    \begin{bmatrix}
      \mathbf{R} & \mathbf{t} \\
      \mathbf{0}^{\mathsf T} & 1
    \end{bmatrix},
    \qquad
    \mathbf{R}\in\mathrm{SO}(3),
    \qquad
    \mathbf{t}\in\mathbb{R}^{3}
  \]
  \begin{itemize}
    \item $\mathbf{R}$은 물체가 어느 방향을 보는지 나타내는 $3 \times 3$ 회전행렬이다.
    \item $\mathbf{t}$는 물체가 어디에 있는지 나타내는 $3 \times 1$ 이동벡터다.
  \end{itemize}
  \vfill
  좌표점 $\mathbf{p}_B$를 좌표계 $A$로 옮기는 식은 다음과 같다.
  \[
    \begin{bmatrix} \mathbf{p}_A \\ 1 \end{bmatrix}
    =
    {}^{A}\mathbf{T}_{B}
    \begin{bmatrix} \mathbf{p}_B \\ 1 \end{bmatrix}
  \]
\end{frame}

\begin{frame}{3. 변환행렬이라는 형식을 쓰는 이유}
  \begin{enumerate}
    \item \textbf{위치와 방향을 하나의 값으로 묶어서 다룬다.}
      따로 관리하면 계산할 때마다 둘을 맞춰줘야 하지만,
      하나로 묶으면 자세 전체를 변수 하나로 취급할 수 있다.
    \item \textbf{좌표계를 갈아타는 일을 곱셈 한 번으로 처리한다.}
      회전만이면 $3 \times 3$ 행렬로 충분하지만, 이동은 곱셈이 아니라 덧셈이라서
      그대로는 섞이지 않는다. 마지막 행에 $1$을 덧붙여 $4 \times 4$로 만들면
      이동까지 곱셈 안으로 들어온다. 그 덕분에 여러 좌표계를 거치는 경로를
      행렬 곱만으로 이어붙일 수 있다.
    \item \textbf{되돌리기가 쉽다.}
      어떤 변환이든 역행렬 하나로 방향을 반대로 뒤집을 수 있다.
  \end{enumerate}
\end{frame}

\begin{frame}{3.1 변환행렬을 곱한다는 뜻}
  경로를 차례대로 연결한다는 뜻이다.
  \[
    {}^{A}\mathbf{T}_{C}
    =
    {}^{A}\mathbf{T}_{B}\,
    {}^{B}\mathbf{T}_{C}
  \]
  \vfill
  이를 행렬 내부까지 펼치면 다음과 같다.
  \[
    \begin{bmatrix}
      \mathbf{R}_1 & \mathbf{t}_1 \\
      \mathbf{0}^{\mathsf T} & 1
    \end{bmatrix}
    \begin{bmatrix}
      \mathbf{R}_2 & \mathbf{t}_2 \\
      \mathbf{0}^{\mathsf T} & 1
    \end{bmatrix}
    =
    \begin{bmatrix}
      \mathbf{R}_1\mathbf{R}_2 & \mathbf{R}_1\mathbf{t}_2+\mathbf{t}_1 \\
      \mathbf{0}^{\mathsf T} & 1
    \end{bmatrix}
  \]
\end{frame}

\begin{frame}{3.2 역행렬의 뜻}
  변환의 방향을 반대로 돌린다.
  \[
    \mathbf{T}^{-1}
    =
    \begin{bmatrix}
      \mathbf{R}^{\mathsf T} & -\mathbf{R}^{\mathsf T}\mathbf{t} \\
      \mathbf{0}^{\mathsf T} & 1
    \end{bmatrix}
  \]
  즉, 회전은 전치행렬 $\mathbf{R}^{\mathsf T}$로 되돌리고,
  이동도 그 회전 좌표계에 맞추어 반대로 적용한다.
  \vfill
  \begin{block}{역행렬을 쓰는 이유}
    \begin{enumerate}
      \item 알고 있는 변환이 필요한 방향과 반대일 때 뒤집어서 쓴다.
      \item 두 자세가 얼마나 다른지 잰다. 자세끼리는 뺄셈이 성립하지 않으므로,
        한쪽에서 다른 쪽까지 얼마나 더 움직여야 하는지를 대신 구한다. (7장)
      \item 여러 변환이 곱해진 식에서 이미 아는 항을 걷어내고 미지수만 남긴다.
        방정식에서 이항하는 것과 같다.
    \end{enumerate}
  \end{block}
\end{frame}

%==================================================================
\section{4. 카메라가 측정하는 것}
%==================================================================

\begin{frame}{4. 카메라가 실제로 측정하는 것은 무엇인가?}
  카메라는 이미지에서 보드나 큐브의 코너를 검출하고 \textbf{PnP}를 풀어서 자세를 얻는다.
  \vfill
  PnP는 Perspective-$n$-Point의 줄임말이다. 한 장의 사진만 가지고 물체가
  카메라 기준으로 어디에 어떤 방향으로 놓여 있는지를 알아내는 방법이다.
  \vfill
  \begin{block}{PnP를 풀려면 두 가지가 필요하다}
    \begin{enumerate}
      \item \textbf{물체 위 점들의 3차원 좌표.} 물체 자신의 좌표계에서 잰 값이다.
        한 변이 $50\,\mathrm{mm}$인 체커보드 격자라면 코너들이
        $(0,0,0)$, $(50,0,0)$, $(0,50,0)$처럼 놓여 있다는 사실을
        설계 도면에서 이미 알고 있다.
      \item \textbf{같은 점들이 사진에서 찍힌 2차원 픽셀 좌표.}
        코너 검출 알고리즘이 찾아준다.
    \end{enumerate}
  \end{block}
\end{frame}

\begin{frame}{4. PnP의 전제와 출력}
  같은 점이 3차원에서는 어디이고 사진에서는 몇 번째 픽셀인지, 그 짝을 여러 개 모으면
  물체의 자세를 역으로 풀 수 있다. 이름의 $n$은 그 짝의 개수를 뜻하며,
  원리상 최소 $3$개, 안정적으로 풀려면 보통 그보다 많이 쓴다.
  \vfill
  단, \textbf{카메라 내부파라미터}를 미리 알고 있어야 한다.
  내부파라미터란 초점거리와 이미지 중심처럼 3차원 점이 이미지의 어느 픽셀에
  맺히는지를 정하는 카메라 자체의 값이다. 이 값은 캘리브레이션 이전에 따로 구해 둔다.
  \vfill
  \begin{alertblock}{정리}
    PnP의 출력이 곧 관측값 $\mathbf{Z}$이다.
    즉 \textbf{카메라 좌표계에서 본} 물체의 자세다.
  \end{alertblock}
\end{frame}

\begin{frame}{4.1 -- 4.3 관측값의 정의}
  \textbf{4.1 고정 카메라 관측}
  \[
    \mathbf{Z}_{i,s} \equiv {}^{C_i}\mathbf{T}_{O,s}
  \]
  세트 $s$의 큐브가 고정 카메라 $i$에서 어떻게 보였는지를 뜻한다.
  \vfill
  \textbf{4.2 그리퍼 카메라 관측}
  \[
    \mathbf{Z}_{g,e} \equiv {}^{C_g}\mathbf{T}_{O,e}
  \]
  이벤트 $e$에서 큐브가 그리퍼 카메라에 어떻게 보였는지를 뜻한다.
  \vfill
  \textbf{4.3 로봇 FK로 알고 있는 그리퍼 자세}
  \[
    \mathbf{G}_e \equiv {}^{B}\mathbf{T}_{G,e}
  \]
  촬영 순간 그리퍼의 위치와 방향이다. 로봇 관절각으로부터 FK를 계산하여 얻는다.
\end{frame}

\begin{frame}{4.3 FK란 무엇인가}
  FK는 Forward Kinematics, 우리말로 \textbf{순기구학}의 줄임말이다.
  \vfill
  로봇의 각 관절이 몇 도 꺾여 있는지를 알 때, 링크 길이 같은 로봇의 설계값을 이용해
  손끝이 베이스 기준으로 어디에 어떤 방향으로 있는지를 계산하는 것이다.
  \vfill
  \begin{block}{중요한 점}
    카메라를 쓰지 않고 \textbf{로봇 자신의 관절 센서만으로} 얻는 값이다.
  \end{block}
\end{frame}

%==================================================================
\section{5. 찾아야 하는 미지수}
%==================================================================

\begin{frame}{5. 우리가 찾아야 하는 미지수}
  \textbf{5.1 고정 카메라 외부파라미터}
  \[
    \mathbf{C}_i \equiv {}^{B}\mathbf{T}_{C_i}
  \]
  \vfill
  \textbf{5.2 그리퍼와 그리퍼 카메라 사이의 hand-eye 변환}
  \[
    \mathbf{X} \equiv {}^{G}\mathbf{T}_{C_g}
  \]
  \vfill
  \textbf{5.3 세트별 큐브 자세}
  \[
    \mathbf{O}_s \equiv {}^{B}\mathbf{T}_{O,s}
  \]
  \vfill
  세 FK 방식의 가장 큰 차이는 $\mathbf{O}_s$를 자유롭게 찾는지,
  FK 값에 고정하는지, 보정된 anchor에 고정하는지다.
\end{frame}

%==================================================================
\section{6. 두 종류 카메라의 관측 경로}
%==================================================================

\begin{frame}{6.1 고정 카메라 경로}
  베이스에서 고정 카메라로 가고, 고정 카메라에서 큐브로 간다.
  \[
    {}^{B}\mathbf{T}_{C_i}\,
    {}^{C_i}\mathbf{T}_{O,s}
    =
    {}^{B}\mathbf{T}_{O,s}
  \]
  앞에서 정의한 짧은 기호를 사용하면 다음과 같다.
  \[
    \mathbf{C}_i\mathbf{Z}_{i,s}=\mathbf{O}_s
  \]
  따라서 고정 카메라가 예측한 큐브 자세는 다음과 같다.
  \[
    \widehat{\mathbf{O}}^{\mathrm{fix}}_{i,s}
    =
    \mathbf{C}_i\mathbf{Z}_{i,s}
  \]
\end{frame}

\begin{frame}{6.2 그리퍼 카메라 경로}
  베이스에서 그리퍼로 가고, 그리퍼에서 카메라로 가고, 카메라에서 큐브로 간다.
  \[
    {}^{B}\mathbf{T}_{G,e}\,
    {}^{G}\mathbf{T}_{C_g}\,
    {}^{C_g}\mathbf{T}_{O,e}
    =
    {}^{B}\mathbf{T}_{O,s(e)}
  \]
  짧은 기호로 쓰면 다음과 같다.
  \[
    \mathbf{G}_e\mathbf{X}\mathbf{Z}_{g,e}
    =
    \mathbf{O}_{s(e)}
  \]
  따라서 그리퍼 카메라가 예측한 큐브 자세는 다음과 같다.
  \[
    \widehat{\mathbf{O}}^{\mathrm{grip}}_e
    =
    \mathbf{G}_e\mathbf{X}\mathbf{Z}_{g,e}
  \]
  \vfill
  \begin{alertblock}{캘리브레이션의 목표}
    같은 세트의 모든 예측은 같은 큐브 자세로 모여야 한다.
  \end{alertblock}
\end{frame}

%==================================================================
\section{7. 두 자세의 차이 재기}
%==================================================================

\begin{frame}{7. 두 자세가 얼마나 다른지 계산하는 방법}
  두 자세 $\mathbf{A}$와 $\mathbf{B}$의 상대변환은 다음과 같다.
  \[
    \mathbf{E} = \mathbf{A}^{-1}\mathbf{B}
  \]
  $\mathbf{A}=\mathbf{B}$라면 $\mathbf{E}=\mathbf{I}$가 된다.
  \vfill
  행렬 내부까지 펼치면 다음과 같다.
  \[
    \mathbf{A}^{-1}\mathbf{B}
    =
    \begin{bmatrix}
      \mathbf{R}_A^{\mathsf T}\mathbf{R}_B
      & \mathbf{R}_A^{\mathsf T}(\mathbf{t}_B-\mathbf{t}_A) \\
      \mathbf{0}^{\mathsf T} & 1
    \end{bmatrix}
  \]
  여기서 회전 부분은 상대 회전이고, 이동 부분은 $\mathbf{A}$ 좌표계에서 본 상대 이동이다.
\end{frame}

\begin{frame}{7.1 이해하기 쉬운 이론식}
  \[
    \boldsymbol{\varepsilon}(\mathbf{A},\mathbf{B})
    =
    \log\!\left(\mathbf{A}^{-1}\mathbf{B}\right)
  \]
  여기서 $\log(\cdot)$는 일반 숫자의 자연로그가 아니다.
  $\mathrm{SE}(3)$ 행렬을 작은 회전 $3$개와 작은 이동 $3$개,
  총 $6$개의 오차 숫자로 바꾸는 Lie logarithm이다.
  \vfill
  개념적으로는 다음과 같이 생각할 수 있다.
  \[
    \log:\mathrm{SE}(3)\rightarrow\mathfrak{se}(3),
    \qquad
    \boldsymbol{\varepsilon}
    =
    \begin{bmatrix} \boldsymbol{\omega} \\ \mathbf{v} \end{bmatrix}
    \in\mathbb{R}^{6}
  \]
  $\boldsymbol{\omega}\in\mathbb{R}^{3}$은 회전오차이고,
  $\mathbf{v}\in\mathbb{R}^{3}$은 이동오차다.
\end{frame}

%==================================================================
\section{8. 공통 목적함수}
%==================================================================

\begin{frame}{8. 모든 방법이 공유하는 기본 목적함수}
  세트 $s$의 공통 큐브 자세를 $\mathbf{Q}_s$라고 하자.
  그러면 vision 관측의 전체 오차는 다음과 같다.
  \[
    \begin{aligned}
      \mathcal{E}_{\mathrm{vis}}
      \left(\{\mathbf{C}_i\},\mathbf{X},\{\mathbf{Q}_s\}\right)
      ={}&
      \sum_{i,s}
      \left\|
        \mathbf{r}\left(\mathbf{C}_i\mathbf{Z}_{i,s},\ \mathbf{Q}_s\right)
      \right\|_2^2
      \\[4pt]
      &+
      \sum_e
      \left\|
        \mathbf{r}\left(\mathbf{G}_e\mathbf{X}\mathbf{Z}_{g,e},\ \mathbf{Q}_{s(e)}\right)
      \right\|_2^2 .
    \end{aligned}
  \]
  \vfill
  아래첨자 vis는 vision, 즉 카메라로 본 정보에서 나온 오차라는 뜻이다.
  로봇 관절값에서 나온 FK 정보와 구분하려고 붙인 이름이다.
\end{frame}

\begin{frame}{8.1 식에 나오는 기호}
  \begin{center}
    \footnotesize
    \begin{tabularx}{\textwidth}{c X l}
      \toprule
      기호 & 뜻 & 값이 어디서 오는가 \\
      \midrule
      $\mathbf{C}_i$ & 고정 카메라 $i$의 자세 (베이스 기준) & 모름. 최적화로 찾아야 함 \\
      $\mathbf{X}$ & 그리퍼와 그리퍼 카메라 사이의 hand-eye 변환 & 모름. 최적화로 찾아야 함 \\
      $\mathbf{Q}_s$ & 세트 $s$의 공통 큐브 자세 (베이스 기준) & FK 방식마다 다름 \\
      $\mathbf{Z}_{i,s}$ & 고정 카메라 $i$가 세트 $s$에서 본 큐브 자세 & 사진에서 PnP를 풀어 얻음 \\
      $\mathbf{Z}_{g,e}$ & 그리퍼 카메라가 이벤트 $e$에서 본 큐브 자세 & 사진에서 PnP를 풀어 얻음 \\
      $\mathbf{G}_e$ & 이벤트 $e$의 그리퍼 자세 (베이스 기준) & 관절각으로 FK를 계산해 얻음 \\
      $\mathbf{r}(\cdot,\cdot)$ & 두 자세의 차이를 $6$개 숫자로 만드는 잔차 함수 & 7장에서 정의한 함수 \\
      $\|\cdot\|_2^2$ & 그 $6$개 숫자를 제곱해서 더한 값 & 잔차로부터 계산 \\
      $\sum_{i,s}$ & 모든 고정 카메라와 모든 세트에 대해 더한다 & 합 기호 \\
      $\sum_e$ & 모든 그리퍼 카메라 촬영 이벤트에 대해 더한다 & 합 기호 \\
      $s(e)$ & 이벤트 $e$가 속한 세트 번호 & 촬영 기록에 남아 있음 \\
      \bottomrule
    \end{tabularx}
  \end{center}
\end{frame}

\begin{frame}{8.1 식을 읽는 순서}
  안쪽부터 밖으로 나간다.
  \begin{enumerate}
    \item $\mathbf{C}_i\mathbf{Z}_{i,s}$로 고정 카메라들이 예측한 큐브 자세를 만든다.
    \item 그 예측을 공통 큐브 자세 $\mathbf{Q}_s$와 비교하여 잔차 $6$개 숫자를 얻는다.
    \item 제곱해서 더해 오차 하나의 숫자로 만든다.
    \item 모든 카메라와 모든 세트에 대해 이 값을 전부 합친다.
    \item 그리퍼 카메라 쪽도 $\mathbf{G}_e\mathbf{X}\mathbf{Z}_{g,e}$로 같은 과정을 거쳐 합친다.
  \end{enumerate}
\end{frame}

\begin{frame}{8.2 이 식이 뜻하는 것}
  첫 번째 합은 고정 카메라 예측을 공통 큐브 자세에 맞춘다.
  두 번째 합은 그리퍼 카메라 예측을 같은 공통 큐브 자세에 맞춘다.
  \vfill
  \begin{block}{캘리브레이션이란}
    $\mathcal{E}_{\mathrm{vis}}$를 최소로 만드는 $\mathbf{C}_i$와 $\mathbf{X}$를 찾는 것이다.
    측정값 $\mathbf{Z}$와 $\mathbf{G}$는 이미 정해져 있으므로 움직일 수 없고,
    오직 미지수만 움직여서 모든 예측이 한 점에 모이게 만든다.
  \end{block}
  \vfill
  \begin{alertblock}{세 FK 방식의 차이}
    관측식 자체는 다르지 않다. $\mathbf{Q}_s$를 무엇으로 정의하느냐가 다르다.
  \end{alertblock}
\end{frame}

%==================================================================
\section{9. 방법 1: No-FK}
%==================================================================

\begin{frame}{9. 방법 1: No-FK}
  \begin{block}{9.1 쉬운 설명}
    로봇이 알려주는 큐브 자세를 사용하지 않는다.
    카메라 위치, hand-eye, 세트별 큐브 자세를 vision 관측만으로 함께 찾는다.
  \end{block}
  \vfill
  \textbf{9.2 정확한 목적함수}
  \[
    \left\{
      \{\widehat{\mathbf{C}}_i\},\ \widehat{\mathbf{X}},\ \{\widehat{\mathbf{O}}_s\}
    \right\}
    =
    \underset{\{\mathbf{C}_i\},\,\mathbf{X},\,\{\mathbf{O}_s\}}{\operatorname{argmin}}
    \ \mathcal{E}_{\mathrm{vis}}
    \left(\{\mathbf{C}_i\},\ \mathbf{X},\ \{\mathbf{O}_s\}\right)
  \]
  이때 공통 자세는 자유변수다.
  \[
    \mathbf{Q}_s=\mathbf{O}_s
  \]
\end{frame}

\begin{frame}{9.2 비슷해 보이는 세 기호의 구분}
  \begin{center}
    \begin{tabular}{cl}
      \toprule
      기호 & 무엇인가 \\
      \midrule
      $\mathbf{Q}_s$ & 카메라 예측을 맞출 기준 자리 \\
      $\mathbf{O}_s$ & 그 자리를 채우는 자유변수 \\
      $\widehat{\mathbf{O}}_s$ & 최적화가 끝난 뒤 나온 추정 결과 \\
      \bottomrule
    \end{tabular}
  \end{center}
  \vfill
  따라서 $\mathbf{Q}_s=\mathbf{O}_s$는 두 값이 같다는 뜻이 아니라,
  기준 자리를 고정값이 아니라 \textbf{자유변수로 채운다}는 선언이다.
  \texttt{Fixed-FK}는 같은 자리를 FK 값으로 채운다.
  \vfill
  셋 중 어느 것도 큐브의 진짜 자세는 아니다.
  진짜 자세는 10.3절에서 $\mathbf{O}^{\mathrm{true}}_s$로 따로 표기한다.
\end{frame}

\begin{frame}{9.3 장점과 주의점}
  \begin{itemize}
    \item raw FK의 계통오차가 캘리브레이션에 직접 들어오지 않는다.
    \item 대신 세트마다 $6$자유도인 $\mathbf{O}_s$가 추가되어 미지수가 많아진다.
    \item 카메라 연결이나 로봇 움직임이 충분하지 않으면 전체 좌표계의
      gauge가 흔들릴 수 있다.
    \item \texttt{No-FK}도 $\mathbf{G}_e$는 사용한다.
      큐브 FK prior $\mathbf{F}_s$만 사용하지 않는다.
  \end{itemize}
\end{frame}

%==================================================================
\section{10. 방법 2: Fixed-FK}
%==================================================================

\begin{frame}{10. 방법 2: Fixed-FK}
  \begin{block}{10.1 쉬운 설명}
    로봇 FK가 알려준 큐브 자세를 정답으로 간주하고 움직이지 못하게 고정한다.
  \end{block}
  \vfill
  raw FK 큐브 prior를 다음과 같이 정의한다.
  \[
    \mathbf{F}_s \equiv {}^{B}\mathbf{T}^{\mathrm{FK}}_{O,s}
  \]
  큐브 자세는 다음 제약을 만족해야 한다.
  \[
    \mathbf{O}_s=\mathbf{F}_s
  \]
\end{frame}

\begin{frame}{10.2 정확한 목적함수}
  \[
    \left\{\{\widehat{\mathbf{C}}_i\},\ \widehat{\mathbf{X}}\right\}
    =
    \underset{\{\mathbf{C}_i\},\,\mathbf{X}}{\operatorname{argmin}}
    \ \mathcal{E}_{\mathrm{vis}}
    \left(\{\mathbf{C}_i\},\ \mathbf{X},\ \{\mathbf{F}_s\}\right)
  \]
  \vfill
  이를 완전히 펼치면 다음과 같다.
  \[
    \begin{aligned}
      \min_{\{\mathbf{C}_i\},\,\mathbf{X}}\ {}&
      \sum_{i,s}
      \left\|\mathbf{r}(\mathbf{C}_i\mathbf{Z}_{i,s},\ \mathbf{F}_s)\right\|_2^2
      \\[4pt]
      &+
      \sum_e
      \left\|\mathbf{r}(\mathbf{G}_e\mathbf{X}\mathbf{Z}_{g,e},\ \mathbf{F}_{s(e)})\right\|_2^2 .
    \end{aligned}
  \]
  \vfill
  \alert{argmin 아래에 $\mathbf{O}_s$가 없다.} 큐브 자세는 FK 값에 못 박혀
  움직일 수 없으므로 후보 목록에서 빠진다.
\end{frame}

\begin{frame}{10.3 왜 FK 오차가 카메라로 전파되는가?}
  실제 큐브 자세를 $\mathbf{O}^{\mathrm{true}}_s$라고 하자.
  raw FK에 계통오차 $\mathbf{D}$가 있으면 다음처럼 쓸 수 있다.
  \[
    \mathbf{F}_s = \mathbf{O}^{\mathrm{true}}_s\mathbf{D}
  \]
  \vfill
  그런데 최적화는 $\mathbf{F}_s$를 움직일 수 없다.
  따라서 남은 오차를 줄이기 위해 $\mathbf{C}_i$나 $\mathbf{X}$가 잘못 움직일 수 있다.
  \vfill
  \begin{alertblock}{맞교환}
    Fixed-FK는 FK가 정확할 때 미지수가 적고 안정적이지만,
    FK에 공통적인 오정렬이 있으면 그 오차를 캘리브레이션 결과에
    흡수시킬 위험이 있다.
  \end{alertblock}
\end{frame}

\begin{frame}{10.4 카메라에도 계통오차가 있으면 (1) 설정}
  카메라도 내부파라미터나 왜곡 보정이 조금 틀리면 항상 같은 방향으로 치우친
  $\mathbf{Z}$를 내놓는다. 예를 들어 큐브를 늘 $3\,\mathrm{mm}$ 멀리 있다고 보고하는 식이다.
  \vfill
  무슨 일이 벌어지는지 숫자로 따라가 보자. 설명을 위해 $1$차원으로 단순화한다.
  \vfill
  \begin{columns}[T]
    \begin{column}{0.48\textwidth}
      \begin{block}{참값}
        \begin{itemize}
          \item 큐브의 진짜 위치: $100\,\mathrm{mm}$
          \item 카메라의 진짜 위치: $0\,\mathrm{mm}$
        \end{itemize}
        카메라가 정직하다면 $\mathbf{Z}=100$이라고 보고해야 한다.
      \end{block}
    \end{column}
    \begin{column}{0.48\textwidth}
      \begin{block}{두 가지 계통오차}
        \begin{itemize}
          \item FK 오차: 로봇은 큐브가 $105$에 있다고 말한다. $5\,\mathrm{mm}$ 초과.
          \item 카메라 오차: 카메라는 큐브가 $103$ 거리에 있다고 본다. $3\,\mathrm{mm}$ 초과.
        \end{itemize}
      \end{block}
    \end{column}
  \end{columns}
\end{frame}

\begin{frame}{10.4 카메라에도 계통오차가 있으면 (2) 계산}
  $\mathbf{Q}_s$는 FK 값 $105$에 못 박혀 움직이지 않는다.
  따라서 $\mathbf{C}\mathbf{Z}=\mathbf{Q}_s$를 만족하도록 카메라 위치를 정한다.
  \[
    \mathbf{C}+103=105
    \quad\Longrightarrow\quad
    \mathbf{C}=2
  \]
  카메라의 진짜 위치는 $0$이었으므로 결과는 $2\,\mathrm{mm}$ 틀렸다.
  \vfill
  \begin{alertblock}{핵심: 최적화가 본 정보의 양}
    최적화가 아는 것은 $105$와 $103$의 차이인 $2$라는 숫자 하나뿐이다.
    그런데 그 $2$의 정체는 다음과 같다.
    \[
      2 = \underbrace{5}_{\text{FK가 초과한 몫}} - \underbrace{3}_{\text{카메라가 초과한 몫}}
    \]
  \end{alertblock}
\end{frame}

\begin{frame}{10.4 카메라에도 계통오차가 있으면 (3) 분리 불가능}
  최적화는 이 분해를 볼 수 없다. $5$와 $3$이든 $7$과 $5$든 $2$와 $0$이든
  모두 같은 $2$를 만들고, 목적함수 안에는 셋을 구별할 정보가 없다.
  만약 카메라 오차가 반대 방향이었다면 $\mathbf{C}=8$이 되어 훨씬 나빠졌을 것이다.
  \vfill
  \begin{block}{구조적으로 분리가 불가능한 이유}
    카메라의 계통오차가 카메라 좌표계에서의 일정한 어긋남이라면,
    그것은 카메라가 다른 곳에 있다고 말하는 것과 수학적으로 같다.
    즉 $\mathbf{C}_i$를 바꾸는 것과 구별되지 않는다.
  \end{block}
  \vfill
  \begin{block}{\texttt{No-FK}와 비교하면}
    \texttt{No-FK}는 큐브 FK 자세를 쓰지 않으므로 오염원이 카메라 하나뿐이다.
    반면 \texttt{Fixed-FK}는 FK 오차와 카메라 오차가 모두 같은 자리인
    $\mathbf{C}_i$와 $\mathbf{X}$로 들어온다.
  \end{block}
\end{frame}

\begin{frame}{10.4 카메라에도 계통오차가 있으면 (4) 잔차의 함정}
  위 계산에서 $\mathbf{C}=2$를 넣으면 $2+103=105$로 정확히 맞아떨어져
  \alert{잔차가 $0$이 된다.}
  수렴은 완벽해 보이지만 카메라 위치는 $2\,\mathrm{mm}$ 틀려 있다.
  \vfill
  \begin{itemize}
    \item 재투영오차는 카메라 자기 자신과의 일치도이므로 이 편향에 둔감하다.
    \item 여러 카메라가 같은 모델과 같은 절차로 캘리브레이션되었다면
      비슷한 방향으로 치우쳐 카메라 간 consistency도 통과한다.
  \end{itemize}
  \vfill
  다만 카메라 오차가 거리나 각도에 따라 달라지는 종류라면 강체 어긋남이 아니므로
  $\mathbf{C}_i$ 하나로 완전히 흡수되지 않고 잔차가 조금 남는다.
  그러나 그 잔차만으로 카메라 탓인지 FK 탓인지 가려내기는 여전히 어렵다.
\end{frame}

%==================================================================
\section{11. 방법 3: Corrected-FK}
%==================================================================

\begin{frame}{11. 방법 3: Corrected-FK}
  \begin{block}{11.1 가장 쉬운 설명}
    raw FK를 무조건 믿지 않는다. 먼저 vision으로 큐브 자세를 계산하고,
    raw FK가 vision과 어떤 공통 차이를 갖는지 학습한다.
    그 차이로 FK를 보정한 뒤, vision과 충분히 가까운 세트에서만
    그 보정된 FK를 사용한다.
  \end{block}
  \vfill
  \begin{block}{전체 흐름}
    \begin{enumerate}
      \item vision-only 초기 해 $\rightarrow$
      \item 세트별 vision 합의 자세 $\mathbf{V}_s$ $\rightarrow$
      \item raw FK와 vision의 공통 차이 $\overline{\boldsymbol{\Delta}}$ $\rightarrow$
      \item FK 보정 $\mathbf{F}^{\mathrm{corr}}_s$ $\rightarrow$
      \item gate 검사 $\rightarrow$
      \item anchor 확정 $\mathbf{A}_s$ $\rightarrow$
      \item anchor 고정 refinement
    \end{enumerate}
  \end{block}
\end{frame}

\begin{frame}{11.2 단계 1: vision-only 초기 해}
  먼저 No-FK 문제를 풀어 초기 카메라와 hand-eye를 구한다.
  \[
    \{\mathbf{C}^{(0)}_i\},\ \mathbf{X}^{(0)},\ \{\mathbf{O}^{(0)}_s\}
    =
    \operatorname{SolveNoFK}(\text{vision observations})
  \]
  \vfill
  위첨자 $(0)$은 초기 추정값이라는 표시다.
  최종 결과가 아니라 중간 단계 값이므로 햇 기호와 구분한다.
\end{frame}

\begin{frame}{11.3 단계 2: 세트별 vision 합의 자세}
  각 세트에서 고정 카메라와 그리퍼 카메라가 예측한 큐브 자세를 모은다.
  \[
    \mathcal{P}_s
    =
    \left\{\mathbf{C}^{(0)}_i\mathbf{Z}_{i,s}\right\}_{i}
    \cup
    \left\{\mathbf{G}_e\mathbf{X}^{(0)}\mathbf{Z}_{g,e}\right\}_{e:\,s(e)=s}
  \]
  \begin{itemize}
    \item $\mathcal{P}_s$는 세트 $s$의 큐브 자세 예측을 출처 구분 없이 전부 담은 주머니다.
      카메라가 $3$대이고 그 세트에서 그리퍼 촬영이 $4$번 있었다면 자세 $7$개가 들어 있다.
    \item $\cup$는 합집합이다. 고정 카메라 예측과 그리퍼 카메라 예측을
      한 통에 담아 동등하게 취급하겠다는 뜻이다.
  \end{itemize}
  \vfill
  이 예측들을 강건 평균하여 vision 합의 자세 $\mathbf{V}_s$를 만든다.
  \[
    \mathbf{V}_s = \operatorname{RobustAverage}\left(\mathcal{P}_s\right)
  \]
  이 단계에서는 예측들을 모두 동등하게 다루고, MAD 기준으로 이상치만 제거한다.
\end{frame}

\begin{frame}{11.3 용어 정리}
  \begin{description}
    \item[이상치] 나머지와 동떨어진 값. 코너 오검출이나 큐브 면 착각으로 생긴다.
      $7$개 중 $6$개가 $2\,\mathrm{mm}$ 안에 모여 있는데 하나가 $100\,\mathrm{mm}$ 튀면
      단순 평균은 약 $14\,\mathrm{mm}$ 밀려난다.
    \item[MAD] Median Absolute Deviation, 중앙값 절대편차.
      값들이 흩어진 정도를 중앙값으로 재기 때문에 표준편차와 달리
      이상치에 오염되지 않는다. 중앙값에서 MAD의 몇 배 이상 떨어진 값을 버린다.
    \item[강건 평균] MAD로 이상치를 거른 뒤 남은 값을 평균한다.
  \end{description}
\end{frame}

\begin{frame}{11.4 단계 3: raw FK와 vision 사이의 공통 차이}
  각 세트의 차이는 다음과 같다.
  \[
    \boldsymbol{\Delta}_s = \mathbf{F}_s^{-1}\mathbf{V}_s
  \]
  즉, raw FK에서 vision 결과로 가려면 얼마나 더 움직여야 하는지를 뜻한다.
  \vfill
  세트별 차이를 강건 가중 평균하여 공통 보정량을 구한다.
  \[
    \overline{\boldsymbol{\Delta}}
    =
    \operatorname{RobustWeightedAverage}_s
    \left(\boldsymbol{\Delta}_s;\ w_s\right)
  \]
  \vfill
  $w_s$는 세트 $s$를 얼마나 믿을지 나타내는 가중치이며,
  그 세트를 뒷받침한 \textbf{관측의 개수}를 쓴다.
  고정 카메라 $3$대가 보고 그리퍼 촬영이 $4$번 있었던 세트라면 $w_s=7$이다.
  카메라 종류는 구분하지 않고 개수만 센다.
\end{frame}

\begin{frame}{11.5 단계 4: FK에 공통 보정량 적용}
  보정량을 오른쪽에 곱한다.
  \[
    \mathbf{F}^{\mathrm{corr}}_s = \mathbf{F}_s\,\overline{\boldsymbol{\Delta}}
  \]
  \vfill
  \begin{block}{세 가지 FK 값의 구분}
    \begin{center}
      \begin{tabular}{cl}
        \toprule
        기호 & 무엇 \\
        \midrule
        $\mathbf{F}_s$ & raw FK. 로봇이 준 원본 \\
        $\mathbf{F}^{\mathrm{corr}}_s$ & 공통 치우침을 걷어낸 FK \\
        $\mathbf{A}_s$ & gate를 거쳐 확정된 최종 anchor \\
        \bottomrule
      \end{tabular}
    \end{center}
  \end{block}
\end{frame}

\begin{frame}{11.6 단계 5: gate로 보정 FK를 믿어도 되는지 검사}
  $\overline{\boldsymbol{\Delta}}$는 모든 세트에 공통으로 적용되는 하나의 보정량이다.
  따라서 세트마다 사정이 다르면 잘 맞지 않을 수 있다.
  큐브가 그리퍼 안에서 미끄러진 세트, 그 세트에서만 관절 오차가 유난히 컸던 경우가 그렇다.
  반대로 그 세트의 관측이 부족해서 vision 쪽이 틀렸을 수도 있다.
  \vfill
  어느 쪽이 틀렸는지는 알 수 없지만, 둘이 크게 어긋났다는 사실은 확인할 수 있다.
  그래서 세트마다 보정된 FK와 vision 합의를 나란히 놓고 차이를 재고,
  그 차이가 기준을 넘으면 해당 세트에서는 FK를 사용하지 않는다.
  이 검사를 \textbf{gate}라고 부른다.
\end{frame}

\begin{frame}{11.6 두 가지 차이를 재는 식}
  이동 차이는 다음과 같다.
  \[
    d_t(s)
    =
    1000
    \left\|
      \mathbf{t}(\mathbf{V}_s) - \mathbf{t}(\mathbf{F}^{\mathrm{corr}}_s)
    \right\|_2
  \]
  \vfill
  회전 차이는 다음과 같다.
  \[
    d_R(s)
    =
    \frac{180}{\pi}
    \cos^{-1}
    \left(
      \frac{
        \operatorname{tr}
        \left(
          \mathbf{R}(\mathbf{V}_s)^{\mathsf T}
          \mathbf{R}(\mathbf{F}^{\mathrm{corr}}_s)
        \right)-1
      }{2}
    \right)
  \]
  \vfill
  세트 $s$는 두 조건을 모두 만족할 때만 통과한다.
  \[
    d_t(s)\leq \tau_t,
    \qquad
    d_R(s)\leq \tau_R
  \]
\end{frame}

\begin{frame}{11.6 기준선은 상수가 아니다}
  기준선 $\tau_t$와 $\tau_R$은 $d_t(s)$와 $d_R(s)$ 값들이
  실제로 어떻게 분포하는지를 보고 정한다.
  \[
    \tau
    =
    \max\left(
      \operatorname{median}_s(d)
      + k \cdot 1.4826 \cdot \operatorname{MAD}_s(d),
      \ \ \tau^{\min}
    \right),
    \qquad
    k=2.5
  \]
  \vfill
  \begin{itemize}
    \item 상수 $1.4826$은 MAD를 표준편차와 같은 척도로 환산하는 값이다.
      정규분포에서 MAD에 이 값을 곱하면 표준편차가 된다.
    \item 이 환산 덕분에 $k$를 표준편차의 배수로 읽을 수 있다.
      $k=2.5$는 중앙값에서 표준편차 $2.5$배까지 정상으로 인정한다는 뜻이며,
      $\overline{\boldsymbol{\Delta}}$를 구하는 강건 평균이 이미 쓰고 있는 값과 같다.
  \end{itemize}
\end{frame}

\begin{frame}{11.6 왜 분포 기준으로 정하는가}
  어떤 세트의 $\boldsymbol{\Delta}_s$가 공통 보정량 $\overline{\boldsymbol{\Delta}}$와
  정확히 같다면 $\mathbf{F}^{\mathrm{corr}}_s=\mathbf{V}_s$가 되어 $d_t(s)=0$이다.
  \vfill
  즉 $d_t(s)$는 FK가 절대적으로 얼마나 틀렸는지가 아니라,
  \textbf{그 세트가 다른 세트들의 공통 경향에서 얼마나 벗어났는지}를 재는 값이다.
  \vfill
  \begin{block}{그렇다면 얼마부터 유별난 값인가}
    나머지 세트들이 얼마나 모여 있느냐에 따라 달라진다.
    \begin{itemize}
      \item 세트들이 모두 $2\,\mathrm{mm}$ 안에 모여 있다면
        $10\,\mathrm{mm}$짜리 세트는 명백히 이상하다.
      \item 원래부터 $30\,\mathrm{mm}$씩 흩어져 있었다면
        같은 $10\,\mathrm{mm}$도 평범하다.
    \end{itemize}
    고정 상수는 이 두 경우를 구분하지 못한다.
  \end{block}
\end{frame}

\begin{frame}{11.6 하한 $\tau^{\min}$이 필요한 이유}
  기준선을 중앙값 근처로만 잡으면 문제가 생긴다.
  중앙값이란 절반은 그보다 작고 절반은 그보다 크다는 뜻이다.
  \vfill
  따라서 세트들이 모두 잘 맞아서 흩어짐이 거의 없으면,
  기준선이 중앙값에 딱 붙어버리고 \alert{멀쩡한 세트의 절반이 탈락한다.}
  걸러낼 이상치가 없는데도 절반이 FK를 못 쓰게 되는 것이다.
  \vfill
  \begin{alertblock}{하한의 역할}
    무언가를 더 걸러내는 것이 아니라,
    \textbf{구별할 수 없는 차이로 걸러내지 않도록 보장하는 것}이다.
  \end{alertblock}
\end{frame}

\begin{frame}{11.6 하한을 무엇으로 정하는가}
  여기에 다시 상수를 쓰면 처음 문제로 돌아간다.
  대신 \textbf{vision 합의 자신의 흔들림}을 쓴다.
  \[
    \tau^{\min}_t
    =
    \operatorname{median}_s
    \left(\operatorname{scatter}_t(\mathcal{P}_s)\right),
    \qquad
    \tau^{\min}_R
    =
    \operatorname{median}_s
    \left(\operatorname{scatter}_R(\mathcal{P}_s)\right)
  \]
  \vfill
  이렇게 두는 이유는 \textbf{분해능} 때문이다.
  어떤 세트에서 카메라들끼리 $1.5\,\mathrm{mm}$씩 어긋나 있다면,
  $\mathbf{V}_s$ 자체가 그만큼 불확실하다.
  이때 FK가 $\mathbf{V}_s$에서 $1\,\mathrm{mm}$ 벗어나 있다고 해도
  그것이 FK의 잘못인지 vision의 잘못인지 구별할 방법이 없다.
  자의 눈금보다 작은 차이를 재려는 것과 같다.
  \vfill
  구별할 수 없는 차이로 세트를 탈락시키는 것은 근거가 없으므로,
  그 아래는 문제 삼지 않는다.
\end{frame}

\begin{frame}{11.7 단계 6: gate 결과에 따라 anchor 확정}
  세트 $s$의 anchor는 gate 통과 여부로 결정된다.
  \[
    \mathbf{A}_s
    =
    \begin{cases}
      \mathbf{F}^{\mathrm{corr}}_s,
        & d_t(s)\leq\tau_t \ \land\  d_R(s)\leq\tau_R, \\[6pt]
      \mathbf{V}_s, & \text{otherwise}.
    \end{cases}
  \]
  \vfill
  \begin{center}
    \begin{tabular}{ccc}
      \toprule
      gate & anchor & FK 사용 \\
      \midrule
      통과 & $\mathbf{F}^{\mathrm{corr}}_s$ & 그대로 사용 \\
      탈락 & $\mathbf{V}_s$ & 사용하지 않음 \\
      \bottomrule
    \end{tabular}
  \end{center}
  \vfill
  통과한 세트는 보정된 FK를 \textbf{그대로} anchor로 삼는다.
  몇 퍼센트만 반영하는 식의 부분 신뢰는 두지 않는다.
  gate가 이미 그 세트를 믿을지 말지 판정했으므로,
  통과한 뒤에 다시 비중을 깎을 근거가 없기 때문이다.
\end{frame}

\begin{frame}{11.8 단계 7: anchor를 고정하고 최종 refinement}
  \texttt{Corrected-FK}는 만들어진 $\mathbf{A}_s$를 최종 단계에서 고정 anchor로 사용한다.
  \[
    \left\{\{\widehat{\mathbf{C}}_i\},\ \widehat{\mathbf{X}}\right\}
    =
    \underset{\{\mathbf{C}_i\},\,\mathbf{X}}{\operatorname{argmin}}
    \ \mathcal{E}_{\mathrm{vis}}
    \left(\{\mathbf{C}_i\},\ \mathbf{X},\ \{\mathbf{A}_s\}\right)
  \]
  \vfill
  \begin{alertblock}{정리}
    Corrected-FK는 단순히 목적함수에 약한 벌점 하나를 더하는 것과 다르다.
    vision 초기 해, FK 공통 정렬, gate 판정, anchor 확정,
    고정 anchor refinement의 순서로 동작한다.
  \end{alertblock}
\end{frame}

\begin{frame}{11.9 아직 검증이 남은 부분}
  이 장의 절차는 확정된 것이 아니라 검증과 보완이 진행 중이다.
  특히 다음 항목들은 실험으로 근거를 마련해야 한다.
  \begin{itemize}
    \item \textbf{$k=2.5$의 타당성.} 현재는 강건 평균이 쓰는 값과 통일했을 뿐이다.
      $k$를 바꿔가며 결과가 얼마나 달라지는지 확인해야 한다.
    \item \textbf{gate가 실제로 도움이 되는가.} 세트를 탈락시키는 것이 결과를
      개선하는지, 아니면 정보를 버리는 손해가 더 큰지는 gate를 끈 조건과
      비교해야 알 수 있다.
    \item \textbf{하한을 vision 산포로 두는 방식의 적절성.}
      산포의 중앙값을 쓰는 것이 맞는지, 세트별 산포를 각각 쓰는 편이 나은지는
      아직 비교하지 않았다.
    \item \textbf{실데이터에서의 확인.} 적응형 기준선은 시뮬레이션과 실제
      캘리브레이션 코드 양쪽에 반영되어 있으나, 실제 촬영 데이터에서
      세트가 어떻게 걸러지는지는 아직 확인하지 않았다.
  \end{itemize}
  \vfill
  따라서 이 장의 수치와 정책은 이후 실험 결과에 따라 바뀔 수 있다.
\end{frame}

%==================================================================

\begin{frame}{한 장 요약: 세 FK 방식}
  모든 방법은 같은 관측식을 쓴다. $\mathbf{Q}_s$를 무엇으로 채우는지가 다를 뿐이다.
  \[
    \mathbf{Q}_s=
    \begin{cases}
      \mathbf{O}_s, & \texttt{No-FK} \quad (\text{자유변수}) \\[4pt]
      \mathbf{F}_s, & \texttt{Fixed-FK} \quad (\text{raw FK에 고정}) \\[4pt]
      \mathbf{A}_s, & \texttt{Corrected-FK} \quad (\text{보정된 anchor에 고정})
    \end{cases}
  \]
  \vfill
  \begin{itemize}
    \item \texttt{No-FK}: 큐브 자세도 직접 찾는다.
    \item \texttt{Fixed-FK}: 큐브 자세를 raw FK에 못 박는다.
    \item \texttt{Corrected-FK}: raw FK를 vision으로 보정하고
      gate를 통과한 세트에서만 사용한다.
  \end{itemize}
\end{frame}

\end{document}
